\Article{20}
Es bleibt nun, aus dem Satz in Artikel 18 die Beschaffenheit der Gleichgewichtsform zu bestimmen. Dabei kommt es vor allem darauf an, allgemein die Variation des Ausdrucks $W$ herzuleiten, wenn die Gestalt des von der Flüssigkeit erfüllten Raumbereichs eine beliebige unendlich kleine Änderung erfährt. Da aber die Variationsrechnung für Doppelintegrale bisher wenig entwickelt ist, wenn auch die Grenzen als veränderlich anzusehen sind, müssen wir dieser schwierigen Frage etwas tiefer nachgehen.

Von der Oberfläche, die den Raumbereich $s$ vom übrigen Raum trennt, betrachten wir den Teil $U$ und setzen voraus, dass jeder seiner Punkte durch drei Koordinaten $x$, $y$, $z$ bestimmt ist, von denen die dritte den Abstand von einer beliebigen horizontalen Ebene angibt. Somit lässt sich $z$ als Funktion der unbestimmten Größen $x$, $y$ auffassen. Ihre partiellen Differenziale bezeichnen wir in der üblichen Weise, jedoch ohne die Klammern, mit
\[
\frac{\mathrm{d}z}{\mathrm{d}x}.\mathrm{d}x,\qquad\frac{\mathrm{d}z}{\mathrm{d}y}.\mathrm{d}y
\]
In jedem Punkt der Oberfläche denken wir uns eine Gerade, die senkrecht zur Oberfläche steht und in Bezug auf den Raumbereich $s$ nach außen gerichtet ist. Die Kosinus der Winkel zwischen dieser Normalen und den zu den Koordinatenachsen $x$, $y$, $z$ parallelen Geraden bezeichnen wir mit $\xi$, $\eta$, $\zeta$. Damit gilt
\[
\begin{gathered}
\xi\xi+\eta\eta+\zeta\zeta=1\\
\frac{\mathrm{d}z}{\mathrm{d}x}=-\frac{\xi}{\zeta},\qquad\frac{\mathrm{d}z}{\mathrm{d}y}=-\frac{\eta}{\zeta}
\end{gathered}
\]
Die Grenze der Oberfläche $U$ ist eine in sich geschlossene Linie, die wir mit $P$ bezeichnen. Wir denken sie uns in einer stetigen Bewegung \SourceLeaf{62}durchlaufen und nehmen ihre Elemente $\mathrm{d}P$ ebenso wie die Oberflächenelemente $\mathrm{d}U$ stets positiv. Die Kosinus der Winkel, welche die Richtung des Elements $\mathrm{d}P$ mit den Koordinatenachsen $x,y,z$ bildet, bezeichnen wir mit $X,Y,Z$. Damit der Sinn dieser Richtung eindeutig ist, legen wir ihn so fest, dass die Richtung selbst an erster Stelle steht. An zweiter Stelle steht die zum Element $\mathrm{d}P$ senkrechte Richtung, die die Oberfläche $U$ berührt und in Bezug auf sie nach innen weist. An dritter Stelle steht schließlich die Normale zur Oberfläche, die in Bezug auf den Raumbereich $s$ nach außen weist. In dieser Reihenfolge bilden die drei Geraden ein System mit derselben Orientierung wie die Koordinatenachsen $x,y,z$. Man sieht leicht (vgl. \textit{Disquiss. gen. circa superficies curvas} Art. 2), dass die Kosinus der Winkel zwischen jener zweiten Richtung und den Koordinatenachsen $x,y,z$ der Reihe nach
\[
\eta^0Z-\zeta^0Y,\qquad \zeta^0X-\xi^0Z,\qquad \xi^0Y-\eta^0X
\]
sind, wenn $\xi^0,\eta^0,\zeta^0$ die Werte von $\xi,\eta,\zeta$ am Punkt des Elements $\mathrm{d}P$ sind.

\Article{21}
Nach diesen Vorbereitungen nehmen wir an, dass die Oberfläche $U$ eine beliebige unendlich kleine Änderung erfährt. Genügte es, nur solche Änderungen zu betrachten, bei denen die Grenze $P$ stets unverändert bliebe oder wenigstens auf derselben senkrechten Fläche läge, so brauchten wir offenbar nur die dritte Koordinate $z$ zu variieren, wodurch die Aufgabe viel einfacher würde. Da wir sie aber in größtmöglicher Allgemeinheit behandeln müssen, würde es bei einem solchen Vorgehen zu umständlichen Verwicklungen kommen, die der Eleganz der Untersuchung schadeten, sobald wir die Veränderlichkeit der Grenzen berücksichtigen. Daher ist es besser, von Anfang an alle drei Koordinaten zu variieren. Wir stellen uns also vor, dass wir jeden Punkt der Oberfläche mit den Koordinaten $x,y,z$ durch einen anderen mit den Koordinaten $x+\delta x$, $y+\delta y$, $z+\delta z$ ersetzen. Dabei lassen sich $\delta x,\delta y,\delta z$ als unbestimmte Funktionen von $x,y$ auffassen, deren Werte jedoch unendlich klein bleiben. Untersuchen wir nun die Variationen der einzelnen Elemente des Ausdrucks $W$, beginnend mit der Variation des Elements $\mathrm{d}U$.

Denken wir uns ein dreieckiges Element der Oberfläche $U$, bezeichnet mit $\mathrm{d}U$, zwischen den Punkten mit den Koordinaten
\[
\begin{array}{lll}
x,&y,&z\\
x+\mathrm{d}x,&y+\mathrm{d}y.&z+\frac{\mathrm{d}z}{\mathrm{d}x}.\mathrm{d}x+\frac{\mathrm{d}z}{\mathrm{d}y}.\mathrm{d}y\\
x+\mathrm{d}'x,&y+\mathrm{d}'y,&z+\frac{\mathrm{d}z}{\mathrm{d}x}.\mathrm{d}'x+\frac{\mathrm{d}z}{\mathrm{d}y}.\mathrm{d}'y
\end{array}
\]
\SourceLeaf{63}Nach bekannten Grundsätzen ergibt sich für den doppelten Flächeninhalt dieses Dreiecks
\[
=(\mathrm{d}x.\mathrm{d}'y-\mathrm{d}y.\mathrm{d}'x)\BareRoot{[1+(\frac{\mathrm{d}z}{\mathrm{d}x})^2+(\frac{\mathrm{d}z}{\mathrm{d}y})^2]}
\]
wenn wir, was zulässig ist, voraussetzen, dass $\mathrm{d}x.\mathrm{d}'y-\mathrm{d}y.\mathrm{d}'x$ eine positive Größe ist.

Auf der variierten Oberfläche haben wir anstelle jener Punkte drei andere. Für ihre Koordinaten erhalten wir
\[
\begin{array}{ll}
\text{des ersten Punktes}&x+\delta x,\quad y+\delta y,\quad z+\delta z
\end{array}
\]
des zweiten Punktes
\[
\begin{gathered}
x+\mathrm{d}x+\delta x+\frac{\mathrm{d}\delta x}{\mathrm{d}x}.\mathrm{d}x+\frac{\mathrm{d}\delta x}{\mathrm{d}y}.\mathrm{d}y\\
y+\mathrm{d}y+\delta y+\frac{\mathrm{d}\delta y}{\mathrm{d}x}.\mathrm{d}x+\frac{\mathrm{d}\delta y}{\mathrm{d}y}.\mathrm{d}y\\
z+\frac{\mathrm{d}z}{\mathrm{d}x}.\mathrm{d}x+\frac{\mathrm{d}z}{\mathrm{d}y}.\mathrm{d}y+\delta z+\frac{\mathrm{d}\delta z}{\mathrm{d}x}.\mathrm{d}x+\frac{\mathrm{d}\delta z}{\mathrm{d}y}.\mathrm{d}y
\end{gathered}
\]
des dritten Punktes
\[
\begin{gathered}
x+\mathrm{d}'x+\delta x+\frac{\mathrm{d}\delta x}{\mathrm{d}x}.\mathrm{d}'x+\frac{\mathrm{d}\delta x}{\mathrm{d}y}.\mathrm{d}'y\\
y+\mathrm{d}'y+\delta y+\frac{\mathrm{d}\delta y}{\mathrm{d}x}.\mathrm{d}'x+\frac{\mathrm{d}\delta y}{\mathrm{d}y}.\mathrm{d}'y\\
z+\frac{\mathrm{d}z}{\mathrm{d}x}.\mathrm{d}'x+\frac{\mathrm{d}z}{\mathrm{d}y}.\mathrm{d}'y+\delta z+\frac{\mathrm{d}\delta z}{\mathrm{d}x}.\mathrm{d}'x+\frac{\mathrm{d}\delta z}{\mathrm{d}y}.\mathrm{d}'y
\end{gathered}
\]
Für den doppelten Flächeninhalt des Dreiecks zwischen diesen Punkten ergibt sich nach derselben Methode
\[
=(\mathrm{d}x.\mathrm{d}'y-\mathrm{d}y.\mathrm{d}'x)\BareRoot{N}
\]
wenn der Kürze halber mit $N$ die Summe
\[
\begin{aligned}
&[(1+\frac{\mathrm{d}\delta x}{\mathrm{d}x})(1+\frac{\mathrm{d}\delta y}{\mathrm{d}y})-\frac{\mathrm{d}\delta x}{\mathrm{d}y}.\frac{\mathrm{d}\delta y}{\mathrm{d}x}]^2\\
+&[(1+\frac{\mathrm{d}\delta x}{\mathrm{d}x})(\frac{\mathrm{d}z}{\mathrm{d}y}+\frac{\mathrm{d}\delta z}{\mathrm{d}y})-\frac{\mathrm{d}\delta x}{\mathrm{d}y}(\frac{\mathrm{d}z}{\mathrm{d}x}+\frac{\mathrm{d}\delta z}{\mathrm{d}x})]^2\\
+&[(1+\frac{\mathrm{d}\delta y}{\mathrm{d}y})(\frac{\mathrm{d}z}{\mathrm{d}x}+\frac{\mathrm{d}\delta z}{\mathrm{d}x})-\frac{\mathrm{d}\delta y}{\mathrm{d}x}(\frac{\mathrm{d}x}{\mathrm{d}y}+\frac{\mathrm{d}\delta z}{\mathrm{d}y})]^2
\end{aligned}
\]
bezeichnet wird. Durch Ausmultiplizieren und Weglassen der Größen zweiter Ordnung erhalten wir
\[
\BareRoot{N}=\BareRoot{[1+(\frac{\mathrm{d}z}{\mathrm{d}x})^2+(\frac{\mathrm{d}z}{\mathrm{d}y})^2]}.\left[1+\frac{L}{1+(\frac{\mathrm{d}z}{\mathrm{d}x})^2+(\frac{\mathrm{d}z}{\mathrm{d}y})^2}\right]
\]
wenn der Kürze halber mit $L$ die Summe
\SourceLeaf{64}\[
\begin{aligned}
&\frac{\mathrm{d}\delta x}{\mathrm{d}x}[1+(\frac{\mathrm{d}z}{\mathrm{d}y})^2]-\frac{\mathrm{d}\delta x}{\mathrm{d}y}.\frac{\mathrm{d}z}{\mathrm{d}x}.\frac{\mathrm{d}z}{\mathrm{d}y}-\frac{\mathrm{d}\delta y}{\mathrm{d}x}.\frac{\mathrm{d}z}{\mathrm{d}x}.\frac{\mathrm{d}z}{\mathrm{d}y}\\
&+\frac{\mathrm{d}\delta y}{\mathrm{d}y}[1+(\frac{\mathrm{d}z}{\mathrm{d}x})^2]+\frac{\mathrm{d}\delta z}{\mathrm{d}x}.\frac{\mathrm{d}z}{\mathrm{d}x}+\frac{\mathrm{d}\delta z}{\mathrm{d}y}.\frac{\mathrm{d}z}{\mathrm{d}y}
\end{aligned}
\]
bezeichnet wird. Das erste Dreieck verhält sich also zum zweiten wie $1$ zu
\[
1+\frac{L}{1+(\frac{\mathrm{d}z}{\mathrm{d}x})^2+(\frac{\mathrm{d}z}{\mathrm{d}y})^2}
\]
und dieses Verhältnis ist daher von der Gestalt des Dreiecks $\mathrm{d}U$ unabhängig. Damit ergibt sich
\[
\delta\mathrm{d}U=\frac{L\mathrm{d}U}{1+(\frac{\mathrm{d}z}{\mathrm{d}x})^2+(\frac{\mathrm{d}z}{\mathrm{d}y})^2}
\]
oder in ausgeschriebener Form
\[
\begin{aligned}
\delta\mathrm{d}U=\mathrm{d}U\big(&\frac{\mathrm{d}\delta x}{\mathrm{d}x}(\eta\eta+\zeta\zeta)-\frac{\mathrm{d}\delta x}{\mathrm{d}y}.\xi\eta-\frac{\mathrm{d}\delta y}{\mathrm{d}x}.\xi\eta\\
&+\frac{\mathrm{d}\delta y}{\mathrm{d}y}(\xi\xi+\zeta\zeta)-\frac{\mathrm{d}\delta z}{\mathrm{d}x}.\xi\zeta-\frac{\mathrm{d}\delta z}{\mathrm{d}y}.\eta\zeta\big)
\end{aligned}
\]

